% Vietnamese translation for OI Public Library
% Source-PDF: IOI中国国家候选队论文/2009/张昆玮/数学归纳法与解题之道.pdf
\documentclass[11pt]{article}
\usepackage{oipl-vn}

\newcommand{\Oh}{\mathcal{O}}
\newcommand{\NP}{\mathsf{NP}}
\newcommand{\True}{\mathsf{True}}
\newcommand{\xor}{\mathbin{\mathrm{xor}}}

\title{Quy nạp toán học và con đường giải bài}
\author{Zhang Kunwei\\Trường Thực nghiệm Sơn Tây\\Huấn luyện viên: Tang Wenbin, Hu Weidong}
\date{16 tháng 12 năm 2008}

\begin{document}
\maketitle

\origtitle{Quy nạp toán học và con đường giải bài}
\sourcepath{IOI China candidate team papers / 2009 / Zhang Kunwei}

\begin{abstract}
Bài viết tóm tắt lý thuyết liên quan đến quy nạp toán học, rồi kết hợp kinh
nghiệm giải bài của tác giả để trình bày ứng dụng của quy nạp trong Olympic
Tin học. Qua nhiều bài toán kinh điển và bài thi, bài viết bàn về chứng minh
tính đúng đắn của thuật toán, thiết kế thuật toán xây dựng, quan hệ giữa quy
nạp và tối ưu hóa thuật toán, cũng như phạm vi áp dụng và các mở rộng của lối
tư duy quy nạp.
\end{abstract}

\tableofcontents

\section{Dẫn nhập}

Trong khi học thuật toán, ta thường tự hỏi: những thuật toán tự nhiên và tinh
tế ấy được nghĩ ra như thế nào, và vì sao chúng đúng? Đằng sau mỗi lời giải
không chỉ có kỹ thuật cụ thể, mà còn có một cách nhìn. Nếu muốn tìm một điểm
khởi đầu cho ``đạo'' giải bài, quy nạp toán học là một lựa chọn rất đáng tin.

Quy nạp không chỉ dùng để viết phần chứng minh sau khi đã biết đáp án. Trong
quá trình đoán, thử, tách bài toán và xây dựng thuật toán, ta cũng nhiều lần
dùng tư tưởng: giải được trường hợp nhỏ, rồi tìm cách mở rộng có kiểm soát
sang trường hợp lớn hơn.

\section{Về quy nạp toán học}

\subsection{Một hồi tưởng ngắn}

Nguyên lý quy nạp toán học có thể truy ngược tới các công trình số học thời
Phục hưng và sau này được đặt trên nền hệ tiên đề Peano. Vì xử lý trực tiếp
những mệnh đề phụ thuộc vào số nguyên dương, quy nạp xuất hiện trong số học,
logic, tổ hợp, đồ thị và nhiều cấu trúc rời rạc khác. Trong Tin học, nơi thuật
toán thường dựa trên đệ quy, cấu trúc dữ liệu và biến đổi kích thước bài toán,
quy nạp càng có vai trò tự nhiên.

\subsection{Các định lý, khái niệm và phương pháp cơ bản}

\paragraph{Quy nạp toán học thứ nhất.}
Với một mệnh đề \(P(n)\) trên các số tự nhiên dương, nếu
\[
  P(1)=\True,\qquad P(n)\Rightarrow P(n+1)\quad(n\in\mathbb{N}^{+}),
\]
thì \(P(n)\) đúng với mọi \(n\). Có thể hình dung các số tự nhiên là một dãy
quân domino: quân đầu tiên ngã, và mỗi quân ngã kéo theo quân kế tiếp.

\paragraph{Quy nạp toán học thứ hai.}
Nếu với mọi \(n\),
\[
  \bigl(\forall k<n,\; P(k)\bigr)\Rightarrow P(n),
\]
thì \(P(n)\) đúng với mọi \(n\). Dạng này linh hoạt hơn vì ở bước \(n\) ta
có thể dùng mọi kết quả đã chứng minh trước đó, không chỉ trường hợp \(n-1\).
Nó cho phép giảm kích thước bài toán đi một nửa, chia thành nhiều phần nhỏ,
hoặc chọn một trường hợp con thuận lợi.

\paragraph{Quy nạp cấu trúc.}
Quy nạp cấu trúc là biến thể thường gặp trong logic, khoa học máy tính và lý
thuyết đồ thị. Ta không nhất thiết quy nạp theo một số tuyến tính, mà theo một
thứ tự topo của cấu trúc, thường là kích thước bài toán. Với cây, ta có thể
tách các cây con; với đồ thị, ta có thể bỏ một cạnh hoặc một đỉnh; với biểu
thức, ta xét các biểu thức con. Tính đúng đắn dựa trên nguyên lý phản ví dụ
nhỏ nhất: nếu có phản ví dụ thì tồn tại phản ví dụ nhỏ nhất, nhưng lập luận
quy nạp lại tạo ra một phản ví dụ nhỏ hơn, mâu thuẫn.

\paragraph{Ví dụ 1: mã Prüfer cho cây gán nhãn.}
Công thức Cayley nói rằng có \(n^{n-2}\) cây không gốc khác nhau trên \(n\)
đỉnh gán nhãn. Mã Prüfer cho một song ánh giữa cây và dãy độ dài \(n-2\).
Lập luận quy nạp rất tự nhiên: chọn lá có nhãn nhỏ nhất, ghi nhãn đỉnh kề
với nó vào dãy, xóa lá đó, rồi dùng giả thiết quy nạp cho cây còn lại. Nếu
cố quy nạp trực tiếp trên công thức đếm, phần chứng minh sẽ nặng hơn nhiều;
quy nạp cấu trúc cho ta đúng đối tượng cần thao tác.

\subsection{Không chỉ tổng kết, mà còn khám phá}

Trong giải bài, một phần là đoán và khám phá, phần còn lại là tổng kết và
chứng minh. Quy nạp thường bị xem như công cụ của giai đoạn thứ hai, nhưng
thực ra khi ta tìm cách loại bỏ một phần tử, thu nhỏ dữ liệu, hoặc tăng dần
kích thước nghiệm, ta đã dùng tư duy quy nạp.

\paragraph{Ví dụ 2: bài toán người nổi tiếng.}
Cho đồ thị có hướng dưới dạng ma trận kề, cần tìm đỉnh chỉ có cạnh đi vào và
không có cạnh đi ra, với số câu hỏi ít nhất. Xác định trực tiếp ứng viên có
thể tốn cả ma trận, nhưng loại một đỉnh không thể là đáp án thì dễ hơn. Hỏi
một cạnh giữa hai ứng viên: nếu cạnh tồn tại, loại đỉnh xuất phát; nếu không,
loại đỉnh kết thúc. Mỗi câu hỏi giảm kích thước bài toán đi \(1\), rồi giả
thiết quy nạp tiếp tục xử lý tập còn lại. Thuật toán tương tác chuẩn này chính
là một quá trình quy nạp.

\section{Ứng dụng trong chứng minh tính đúng đắn}

Trong luyện tập ta thường lấy AC làm điểm kết thúc, nhưng hiểu vì sao thuật
toán đúng mới là nền tảng vững chắc. Đặc biệt với tham lam, chứng minh thường
quyết định thuật toán có đáng tin hay không. Quy nạp là một trong các công cụ
chứng minh phổ biến nhất của thuật toán.

\subsection{Thuật toán tham lam}

\paragraph{Ví dụ 3: Robbers.}
\(N\) tên cướp chia \(M\) đồng vàng theo tỉ lệ \(X_i/\sum X_i\), nhưng vàng
không thể chia nhỏ. Gọi \(K_i\) là số vàng tên thứ \(i\) nhận được; cần làm
nhỏ sai lệch tổng thể. Một chiến lược tự nhiên là trước hết chia phần nguyên
theo tỉ lệ, rồi mỗi lần cấp một đồng vàng cho người có độ thiếu hụt
\[
  \Delta_i=\frac{X_i}{\sum X_i}-\frac{K_i}{\sum K_i}
\]
lớn nhất.

Có thể chứng minh bằng quy nạp rằng sau khi đã phát \(m\) đồng vàng,
\[
  \max_i\Delta_i-\min_i\Delta_i\le \frac1M.
\]
Ở cơ sở, việc chia tỉ lệ ban đầu thỏa bất đẳng thức. Khi thêm một đồng cho
người đang thiếu nhất, nếu người đó chưa vượt lên thành người ``dư'' nhất thì
khoảng cách chỉ giảm. Nếu có vượt, giá trị mới cũng chỉ tăng thêm đúng phần
đóng góp của một đồng vàng, nên bất đẳng thức vẫn giữ được. Phần còn lại là
chứng minh mọi nghiệm tối ưu cũng phải có tính cân bằng tương tự, từ đó thuật
toán tham lam là đúng.

\subsection{Các thuật toán khác}

\paragraph{Ví dụ 4: Think Positive.}
Một hành tinh có \(n\) ngày trong năm; ngày vui có giá trị \(+1\), ngày buồn
có giá trị \(-1\). Chọn ngày đầu năm sao cho mọi tổng tiền tố trong năm đều
dương. Bài toán trở thành đếm số vị trí bắt đầu trong một dãy vòng sao cho
tổng tiền tố luôn dương.

Xét quy nạp mạnh theo độ dài dãy. Nếu tồn tại cặp liên tiếp \(+1,-1\), xóa
cặp này không đổi tổng toàn dãy. Hai vị trí vừa xóa không thể là điểm bắt đầu
hợp lệ, còn các vị trí khác giữ nguyên điều kiện tiền tố hoặc trở nên tốt hơn;
vì vậy bài toán giảm đi \(2\). Nếu không có cặp \(+1,-1\), hoặc dãy toàn
\(-1\) nên đáp án là \(0\), hoặc dãy toàn \(+1\) nên mọi vị trí đều hợp lệ.
Lập luận này cho một thuật toán tuyến tính rất đơn giản.

\section{Ứng dụng trong thuật toán xây dựng}

Thuật toán xây dựng thường khó vì không có khuôn mẫu cố định. Quy nạp giúp ta
tìm điểm bám: xây được nghiệm nhỏ, rồi thêm một phần tử hoặc nối thêm một
khối sao cho bất biến vẫn đúng.

\subsection{Khôi phục cấu trúc dữ liệu}

\paragraph{Ví dụ 5: Set Cover từ bài toán phủ đoạn.}
Bài toán phủ đoạn chọn ít điểm nguyên nhất sao cho mỗi đoạn cho trước chứa
ít nhất một điểm. Nếu một instance của set cover được tạo ra từ bài toán phủ
đoạn, ta muốn khôi phục đủ cấu trúc tuyến tính để dùng tham lam.

Với các đoạn không chứa nhau, duyệt từ trái sang phải và khi gặp đoạn chưa
được phủ thì chọn điểm phải nhất còn phủ được nó. Vì set cover tổng quát là
\(\NP\)-đầy đủ, trọng tâm là khôi phục quan hệ tuyến tính ẩn. Ta chỉ cần biết
thứ tự tương đối của các đoạn và đầu phải của mỗi đoạn. Giả thiết quy nạp có
thể là: trong một thành phần liên thông theo quan hệ giao nhau, ta đã biết vị
trí các đoạn đã xử lý và biết một đoạn \(A\) giao với đoạn mới \(X\). Khoảng
cách giữa \(A\) và \(X\) suy ra từ sai khác đối xứng của hai tập. Để biết
\(X\) nằm phía nào của \(A\), dùng thêm một đoạn đã biết hướng \(B\); \(X\)
và \(B\) ở cùng phía của \(A\) khi phần giao của chúng với \(A\) có quan hệ
bao hàm. Bằng cách tăng cường giả thiết quy nạp, ta lần lượt dựng lại cấu trúc
cần thiết.

\paragraph{Ví dụ 6: Roman Roads.}
Cho quan hệ kề giữa các con đường, cần xác định liệu chúng có thể là các cạnh
của một cây đường La Mã hay không và dựng cây đó. Với cây, quan hệ giữa cạnh
và đỉnh có tính đệ quy rõ rệt. Gọi \(X\) là cạnh nối đỉnh \(x\) với cha \(p\),
và \(P\) là cạnh nối \(p\) với cha của nó. Giả thiết quy nạp: khi biết \(x\),
\(X\) và \(P\), ta có thể dựng cây con gốc \(x\).

Các cạnh kề với \(X\) chia làm hai loại: cạnh cùng nối tại \(p\), và cạnh nối
tại \(x\). Các cạnh loại đầu đã thuộc phần phía trên; các cạnh loại sau tạo
ra các cây con cần tiếp tục quy nạp. Cạnh \(P\) giúp phân biệt hai loại: cạnh
nối tại \(p\) kề với \(P\), còn cạnh nối tại \(x\) thì không. Từ đó có thể
dùng phép hiệu tập để tách nhánh và kiểm tra tính hợp lệ.

\subsection{Xây dựng chiến lược và phương án}

\paragraph{Ví dụ 7: Royal Federation.}
Trong một cây thành phố, cần chia thành các bang, mỗi bang có từ \(B\) đến
\(3B\) thành phố và có một trung tâm hành chính sao cho mọi thành phố trong
bang đi tới trung tâm mà không ra khỏi bang. Hãy xem cây có gốc. Giả thiết
quy nạp: với một cây con nhỏ hơn, ta chia được các phần có kích thước
\(B\le x\le 2B\), còn lại một tập thành phố nối với gốc và có kích thước
\(x\le B\).

Áp dụng giả thiết cho từng cây con của gốc, rồi gom các phần dư. Mỗi khi đủ
\(B\) thành phố thì tạo một bang với trung tâm là gốc hiện tại. Vì mỗi phần
dư con có kích thước không quá \(B\), bang mới không vượt \(2B\); phần dư
cuối cùng không quá \(B\). Ở bước cuối cùng, gộp phần dư vào bang kề nó để
đạt giới hạn \(3B\). Điểm mấu chốt là chọn đúng bất biến ``phần còn lại nối
với gốc và có kích thước không quá \(B\)''.

\paragraph{Ví dụ 8: Electricity.}
Mạng điện là một DAG. Chính phủ muốn đảo hướng một số cạnh, mỗi ngày một
cạnh, và tại mọi thời điểm đồ thị vẫn phải không có chu trình. Giả sử không
có cạnh song song. Quy nạp theo số đỉnh: chọn một đỉnh \(A\) không có cạnh ra
trong trạng thái cuối. Các cạnh đi vào \(A\) không cần đổi, còn các cạnh đi
ra từ \(A\) cần đảo.

Để đảo các cạnh xuất phát từ \(A\), dùng một quy nạp phụ theo số cạnh còn cần
đổi. Chọn cạnh \(A\to B\) có đỉnh cuối \(B\) đứng sớm nhất trong thứ tự topo
của đồ thị ban đầu, rồi đảo thành \(B\to A\). Nếu tạo chu trình, chu trình
phải có dạng
\[
  A\to C\to \cdots \to B\to A.
\]
Nhưng khi đó \(C\) phải đứng sau \(B\) trong thứ tự topo, mâu thuẫn với đường
đi \(C\to\cdots\to B\). Do đó bước đảo này an toàn, và phần còn lại giao cho
giả thiết quy nạp. Từ phân tích này suy ra thuật toán gọn: sắp các cạnh cần
đổi theo thứ tự nghịch của đỉnh đầu trong trạng thái cuối, rồi theo thứ tự
thuận của đỉnh cuối trong trạng thái đầu.

\section{Quy nạp toán học và tối ưu hóa thuật toán}

Tối ưu hóa ở đây không phải tối ưu hằng số bằng hợp ngữ, mà là chọn cách cài
đặt, mô hình và hướng xử lý hiệu quả hơn. Nhiều cải tiến thuật toán có thể
được nhìn dưới lăng kính quy nạp: chọn lại đối tượng quy nạp, cơ sở quy nạp,
hướng quy nạp, hoặc tăng cường giả thiết.

\subsection{Chọn khéo đối tượng quy nạp}

\paragraph{Ví dụ 9: Explaining the Stock Market.}
Cho một dãy giá, cần chia thành ít dãy con đơn điệu tăng hoặc giảm nhất, với
điều kiện một ngày đơn lẻ không được tính là dãy đơn điệu trừ các trường hợp
đặc biệt. Nếu quy nạp theo độ dài còn lại bằng cách mỗi lần lấy ra một dãy
con đơn điệu, ta có một tìm kiếm đúng nhưng nhiều trạng thái.

Trước hết xét biến thể chỉ cho dãy tăng. Duyệt từ trái sang phải; khi thêm
một giá mới, ta đặt nó vào dãy tăng có phần tử cuối lớn nhất nhưng vẫn nhỏ
hơn giá mới. Dãy có đuôi nhỏ hơn còn nhiều tiềm năng hơn, nên không nên dùng
nó nếu không cần. Với cả dãy tăng và giảm, nghiệm cục bộ gồm các đuôi của
các dãy hiện có; mỗi phần tử mới thử ghép vào một dãy tăng hoặc một dãy giảm,
rồi giữ nghiệm có số dãy ít hơn. Như vậy đối tượng quy nạp là tiền tố đã xử
lý và ``nghiệm còn nhiều tiềm năng nhất'', thay vì toàn bộ phần dãy còn lại.

\subsection{Hoàn thiện cơ sở quy nạp}

\paragraph{Ví dụ 10: tối ưu quicksort.}
Quicksort có thể hiểu là thuật toán quy nạp: phân hoạch dãy rồi sắp xếp hai
phần nhỏ hơn. Nhưng với dãy rất ngắn, chi phí đệ quy và phân hoạch có thể lớn
hơn lợi ích. Ta mở rộng cơ sở quy nạp: nếu độ dài nhỏ hơn một ngưỡng, dùng
insertion sort hoặc heap sort. Cơ sở quy nạp xuất hiện rất nhiều lần trong
quá trình chạy, nên đầu tư tối ưu ở đây thường có hiệu quả. Introsort và nhiều
biến thể sắp xếp thực tế dùng đúng tư tưởng này.

\subsection{Chọn hướng quy nạp cẩn thận}

\paragraph{Ví dụ 11: quy tắc Horner.}
Tính đa thức \(\sum_{i=0}^{N} a_i x^i\). Nếu quy nạp theo tổng các hạng từ
bậc thấp lên, ta còn phải duy trì \(x^k\). Đổi hướng từ bậc cao xuống cho
\[
  \sum_{i=0}^{N}a_i x^i
  =(\cdots((a_Nx+a_{N-1})x+a_{N-2})\cdots)x+a_0,
\]
ta được quy tắc Horner, giảm số phép nhân và trạng thái phụ.

\paragraph{Ví dụ 12: xây dựng heap nhị phân.}
Từ một mảng vô thứ tự, xây heap từ trên xuống có độ phức tạp
\(\Theta(n\log n)\), còn xây từ dưới lên là \(\Theta(n)\). Nhìn theo quy nạp,
phần xử lý sớm thường xảy ra khi kích thước bài toán nhỏ, nên có lợi hơn.
Heap có nhiều nút ở đáy; chọn hướng từ dưới lên đặt lợi thế này vào đúng phần
chiếm số lượng lớn.

\subsection{Tăng cường giả thiết quy nạp}

\paragraph{Ví dụ 13: Yet Another Digit.}
Giữ trọng số nhị phân nhưng cho phép mỗi chữ số là \(0,1,2\); số cách biểu
diễn một số \(n\) theo dạng này gọi là \(f(n)\). Công thức trực tiếp là
\[
f(n)=
\begin{cases}
  f(n/2)+f(n/2-1), & n\equiv 0 \pmod 2,\\
  f(\lceil n/2\rceil), & n\equiv 1 \pmod 2.
\end{cases}
\]
Với \(n\) rất lớn, cách này vẫn khó vì các bài toán con chồng lấn và không
thuận tiện cho số nguyên lớn. Quan sát \(f(n)\) thường đi cặp với \(f(n-1)\),
đặt \(g(n)=f(n-1)\) và tăng cường giả thiết quy nạp thành: đã biết
\[
  f(\lceil n/2\rceil),\qquad g(\lceil n/2\rceil).
\]
Khi đó
\[
f(n)=
\begin{cases}
  f(\lceil n/2\rceil)+g(\lceil n/2\rceil), & n\equiv 0 \pmod 2,\\
  f(\lceil n/2\rceil), & n\equiv 1 \pmod 2,
\end{cases}
\]
và
\[
g(n)=
\begin{cases}
  g(\lceil n/2\rceil), & n\equiv 0 \pmod 2,\\
  g(\lceil n/2\rceil)+f(\lceil n/2\rceil), & n\equiv 1 \pmod 2.
\end{cases}
\]
Cài đặt chỉ cần hai biến \(f,g\) và quét biểu diễn nhị phân của số lớn từ cao
đến thấp: gặp bit \(0\) thì cập nhật \(f:=f+g\), gặp bit \(1\) thì cập nhật
\(g:=f+g\). Độ phức tạp trở thành \(\Theta(\log n)\).

\section{Tác dụng gợi mở và giá trị thẩm mỹ}

Các ví dụ trên cố gắng bóc tách những ``nét bút thần'' trong lời giải thành
những bước tự nhiên hơn: chọn cấu trúc nhỏ, xác định bất biến, mở rộng có kiểm
soát. Thuật toán đẹp không nhất thiết là thuật toán gây bất ngờ, mà thường là
thuật toán đơn giản, đi thẳng vào bản chất. Quy nạp là một trong những cách
giúp ta tìm sự đơn giản ấy.

Không thể tồn tại một ``thuật toán tìm mọi thuật toán'' theo nghĩa tuyệt đối,
nhưng ta vẫn có thể rèn luyện các chiến lược tư duy chung. Nếu các định lý và
kỹ thuật cụ thể là con đường trong từng bài, thì quy nạp là một phương pháp
định hướng trên con đường đó.

\section{Vấn đề và hạn chế}

\subsection{Về mặt lý thuyết}

Tính đúng đắn của quy nạp dựa vào tiên đề quy nạp trong hệ Peano. Theo định
lý bất toàn của Gödel, một hệ tiên đề đủ mạnh không thể tự chứng minh toàn bộ
tính đúng đắn của chính nó trong nội bộ hệ. Trong phạm vi thuật toán thi đấu,
ta hầu như không gặp những vấn đề nền tảng kiểu này, nên quy nạp vẫn là công
cụ suy luận đáng tin. Dù vậy, đặc điểm ``từ hữu hạn đi tới vô hạn'' mà không
có chứng minh trực tiếp bên ngoài hệ tiên đề vẫn là một giới hạn triết học
cần biết.

\subsection{Về mặt ứng dụng}

Nếu gặp mọi bài đều viết đầy đủ giả thiết quy nạp, cơ sở và bước chuyển, ta
có thể tự tạo thêm chi phí không cần thiết. Mục tiêu là dùng linh hoạt: khi
cần chứng minh tham lam, dựng thuật toán, hoặc kiểm soát một bất biến khó,
quy nạp phải xuất hiện đúng lúc; khi mô hình đã rõ, ta không cần hình thức hóa
quá mức.

\subsection{Các lớp bài không thuận lợi}

Một số lớp bài không dễ tiếp cận bằng quy nạp, chẳng hạn quy hoạch tuyến tính,
quy hoạch nguyên và các bài quy về chúng như một số mô hình luồng mạng; một
số bài tối ưu hóa toàn cục; phần lớn thuật toán xấp xỉ hoặc không hoàn hảo;
thuật toán số; và các bài dựa chủ yếu trên ma trận hoặc mô hình toán học khác.
Danh sách này không tuyệt đối, vì kinh nghiệm của tác giả có hạn, nhưng nó
nhắc ta rằng quy nạp là một công cụ mạnh chứ không phải chìa khóa vạn năng.

\section{Lời cuối}

Tác giả từng đọc sách nghiên cứu giải bài của thầy Shan Zun, trong đó nhấn
mạnh rằng lời giải nên tự nhiên, đơn giản và đi vào lõi vấn đề. Bài viết ban
đầu chỉ định trình bày một chuyên đề về quy nạp toán học, nhưng trong khi thử
vén màn ``đạo giải bài'', nội dung đã mở rộng thành một thảo luận dài hơn.
Tác giả hy vọng những ví dụ này có ích cho người đọc, đồng thời cảm ơn các
thầy Wang Jinsheng, Ma Hongchun, Tang Wenbin, Hu Weidong và mọi người đã động
viên, góp ý. Mọi trao đổi có thể gửi tới tác giả theo địa chỉ
\texttt{aceeca.135531@gmail.com}.

\section{Nguồn bài tập}

Ngoài các bài kinh điển, một số ví dụ trong bài lấy từ Andrew Stankevich's
Contest. Vì số lượng bài được nhắc tới khá nhiều, bản gốc chỉ giữ lại phần
nội dung cốt lõi liên quan đến chủ đề của bài viết; người đọc quan tâm có thể
tìm đề gốc, lời giải hoặc mã nguồn tương ứng.

\end{document}
